\subsection{No fixed finite odd prime support works uniformly}
\label{sec:no-fixed-prime-support}

\paragraph{Finding A22 (all exponents and all products).}
Let $S$ be any finite set of odd primes and let
$f\in\{1,2\}$.  There is a finite field $\mathbf F_Q$ of
characteristic two and a smooth projective geometrically
connected genus-two curve $C/\mathbf F_Q$ of $2$-rank $f$
such that, for $X=C^{(1)}$, $J=\operatorname{Pic}^0(X)$ and
its Raynaud theta support $\Theta$,
\[
       \Theta(\overline{\mathbf F}_Q)\cap
       \bigl(J(\overline{\mathbf F}_Q)\{S\}\setminus\{0\}\bigr)
                         =\varnothing.                       \tag{SA.1}
\]
Here $J\{S\}$ allows arbitrary powers of the primes in
$S$ and arbitrary products of those powers.
The curve can be chosen in any prescribed nonempty open
of any geometric irreducible component of
$\mathcal M_2^f$.  In particular these examples are
Zariski dense in each such component.

This strengthens finite-list avoidance: no fixed finite
prime support, even with all exponents and mixed orders
allowed, supplies a universal pure odd-torsion theta
witness on either positive-rank stratum.  The assertion
is unconditional.  The field and curve may depend on $S$;
no uniform degree bound for that field is asserted.

\begin{proof}
The case $S=\varnothing$ is immediate, so assume that $S$
is nonempty and put $N=\prod_{\ell\in S}\ell$.

\emph{1. A family with a finite avoidance condition.}
Start with a component of $\mathcal M_2^f$ over
$\overline{\mathbf F}_2$.  Use the tricanonical frame scheme
and choose a nonempty smooth affine open $T$ as in A21.
There is a universal smooth proper genus-two curve
$\mathcal C\to T$, every fibre has $2$-rank $f$, and the
geometric generic monodromy on every odd level is full.
One can also restrict to the inverse image of any
prescribed nonempty moduli open.

Write $\mathcal X=\mathcal C^{(1)}$ for relative Frobenius
twist over $T$ and
$\mathcal J=\operatorname{Pic}^0(\mathcal X/T)$.
Relative Frobenius is finite locally free of degree two.
In the sequence
\[
 0\longrightarrow\mathcal O_{\mathcal X}
  \longrightarrow F_*\mathcal O_{\mathcal C}
  \longrightarrow\mathcal B\longrightarrow0,                 \tag{SA.2}
\]
the quotient $\mathcal B$ is locally free of rank one.
This can be checked in a relative etale coordinate, where
$1,t$ is the usual basis for relative Frobenius.
It has degree one on every fibre.  Consequently
\[
 \mathcal X\longrightarrow\mathcal J,\qquad
 R\longmapsto\mathcal O_{\mathcal X_t}(R)
                         \otimes\mathcal B_t^{-1}             \tag{SA.3}
\]
is a relative translated Abel map.  Its construction uses
the relative diagonal and the Picard functor; a rational
base point on the fibres is unnecessary.  Its image is
closed, because $\mathcal X$ is proper over $T$ and
$\mathcal J$ is separated over $T$, and its geometric
fibre image is precisely the Raynaud theta support.

The scheme
$\mathcal J[N]^\times=\mathcal J[N]\setminus\{0\}$
is finite etale over $T$: the zero section is open and
closed in this finite etale scheme.  The inverse image
of the closed image of (SA.3) in
$\mathcal J[N]^\times$ is closed, and its projection to
$T$ is closed because the projection is finite.
It misses the generic point by A21.  Remove this proper
closed subset and replace what remains by a nonempty
affine open $U$.  Then every geometric fibre over $U$
satisfies
\[
                  \Theta_t\cap
                  (\mathcal J_t[N]\setminus\{0\})
                           =\varnothing.                     \tag{SA.4}
\]
This is a finite incidence condition, so no passage from
infinitely many open subsets to one open has occurred.

All data descend to a sufficiently large finite field
$k_0=\mathbf F_{q_0}$.  Enlarge that constant field so
that $U$ is geometrically integral and
$q_0\equiv1\pmod{N^2}$.  The latter is possible because
$N$ is odd.  The family, its fixed-rank condition,
and (SA.4) are preserved by this constant extension.

\emph{2. A connected symplectic level torsor.}
Choose a primitive $N^2$th root of unity in $k_0$ to
identify the values of the Weil pairing with the standard
constant pairing.  Let
\[
 V=\operatorname{Isom}_{\mathrm{sympl}}
       \bigl((\mathbf Z/N^2)^4_U,\mathcal J[N^2]\bigr)
                      \longrightarrow U.                    \tag{SA.5}
\]
This is a finite etale torsor for the constant group
\[
       G=\operatorname{Sp}_4(\mathbf Z/N^2)
        =\prod_{\ell\in S}
                       \operatorname{Sp}_4(\mathbf Z/\ell^2).
\]
Its geometric generic monodromy is all of $G$ by A21.
The restriction of $V_{\overline{k}_0}$ to the generic
point of $U_{\overline{k}_0}$ is therefore connected.
Since $U$ is smooth and
geometrically integral, $V$ is geometrically integral.
Equivalently, its geometric connectedness follows from
the full monodromy, and smoothness makes distinct
irreducible components disjoint.
The total arithmetic image is also $G$: it contains
the full geometric image, and it is contained in $G$
because $\mu_{N^2}$ is constant over $k_0$.

\emph{3. Realizing a precise Frobenius class by a twist.}
For $\ell\in S$ put $u_\ell=1+\ell$ modulo $\ell^2$ and
choose
\[
 \gamma_\ell=
 \operatorname{diag}(u_\ell,u_\ell,
                         u_\ell^{-1},u_\ell^{-1})
     \in\operatorname{Sp}_4(\mathbf Z/\ell^2).
                                                               \tag{SA.6}
\]
Use the standard symplectic basis ordered as
$e_1,e_2,f_1,f_2$.  Then
\[
 \gamma_\ell\equiv I\pmod\ell,\qquad
 (\gamma_\ell-I)/\ell
                   \equiv\operatorname{diag}(1,1,-1,-1)
                             \pmod\ell.                       \tag{SA.7}
\]
The second matrix is invertible.
Let $\gamma=(\gamma_\ell)_{\ell\in S}\in G$.
Each $\gamma_\ell$ has order $\ell$, so $\gamma$ has
order $N$.

Here is a complete existence argument for a specialization
with this Frobenius class.  It is the finite-field
twisting proof of Chebotarev, with only qualitative
Lang--Weil needed.  Let $\mathcal F$ denote arithmetic
$q_0$-power Frobenius acting on the geometric points of
$V$.  For a symplectic frame $v$ define a new descent
action of the arithmetic generator by
\[
                       v\longmapsto
                               \mathcal F(v)\gamma^{-1}.
                                                               \tag{SA.8}
\]
Right multiplication by $\gamma^{-1}$ is a deck
automorphism defined over $k_0$ and commutes with
$\mathcal F$.  Since $\gamma^N=1$, this descent action
agrees with the original action on
$\operatorname{Gal}(\overline{k}_0/\mathbf F_{q_0^N})$.
It therefore defines a finite Galois descent datum on
$V_{\mathbf F_{q_0^N}}$.  The datum is effective because
$V$ is affine; the exact general reference is the
Stacks Project, Lemma~35.37.1, Tag~\texttt{0245}
\cite{Stacks0245}.
Write $V^\gamma$ for the descended twist.
The projection to $U$ descends because deck
automorphisms act trivially there.

The twist $V^\gamma$ is geometrically integral, being
geometrically isomorphic to $V$.  Lang--Weil gives
$V^\gamma(\mathbf F_{q_0^m})\ne\varnothing$ for every
sufficiently large $m$.  One precise reference for this
form, including affine varieties, is Poonen,
\emph{Rational points on varieties}, Theorem~7.7.1(ii)--(iii),
printed p.~222 \cite{Poonen2017RPV}.  Choose such an $m$
with $m\equiv1\pmod N$ and a rational point of
$V^\gamma$ over $\mathbf F_{q_0^m}$.  Under the geometric
identification with $V$, it gives a frame $v$ over a
point $t\in U(\mathbf F_{q_0^m})$ such that
\[
                \mathcal F^m(v)\gamma^{-m}=v,
       \qquad\text{hence}\qquad
                \mathcal F^m(v)=v\gamma.                      \tag{SA.9}
\]
Thus the $Q$-power Frobenius of the fibre, where
$Q=q_0^m$, has matrix $\gamma$ on its $N^2$-torsion in
the frame $v$.  Changing the frame conjugates this
matrix, which leaves the properties in (SA.7) unchanged.
This argument does not require $t$ to have minimal
residue-field degree $m$; the fibre is considered over
the stated field $\mathbf F_Q$.

\emph{4. All rational $S$-primary torsion has squarefree
exponent.}
Let $C/\mathbf F_Q$ be this fibre and
$\Phi$ its $Q$-power Frobenius on $J$.
For each $\ell\in S$, lift the chosen torsion frame to
a basis of $T_\ell J$.  Equations (SA.7)--(SA.9) imply
\[
             \Phi-I=\ell A_\ell,\qquad
                         A_\ell\in
                         \operatorname{GL}_4(\mathbf Z_\ell).
                                                               \tag{SA.10}
\]
On
$J[\ell^\infty](\overline{\mathbf F}_Q)
 \simeq(\mathbf Q_\ell/\mathbf Z_\ell)^4$,
the matrix $A_\ell$ is an automorphism.  Consequently
\[
 \begin{split}
 J(\mathbf F_Q)[\ell^\infty]
     &=\ker\bigl(\Phi-I\mid
                  J[\ell^\infty](\overline{\mathbf F}_Q)\bigr)\\
     &=J[\ell](\overline{\mathbf F}_Q).
 \end{split}                                                   \tag{SA.11}
\]
Indeed $\ell A_\ell v=0$ if and only if
$\ell v=0$.  In particular all the $\ell$-torsion
points are rational, and there are no rational points
of order $\ell^2$ or higher.  Taking the direct sum
over $S$ gives
\[
                  J(\mathbf F_Q)\{S\}
                          =J[N](\overline{\mathbf F}_Q).
                                                               \tag{SA.12}
\]
This equation includes mixed orders dividing $N$.

\emph{5. Excluding arbitrary exponents geometrically.}
Since $\Phi$ is the identity on $J[\ell]$ for every
$\ell\in S$, Finding~A19 applies with $d=1$.
It puts every geometric $S$-primary theta point in
$J(\mathbf F_Q)\{S\}$.  By (SA.12) this is $J[N]$,
and (SA.4) excludes every nonzero point of that group
from theta.  This proves (SA.1).

The construction started inside an arbitrary prescribed
nonempty open of a component.  Repeating it for each
such open proves the asserted Zariski density.
\end{proof}

\paragraph{What the result leaves open.}
A22 is a nonexistence result for a \emph{fixed finite}
choice of prime factors across all curves.  It does not
exclude a witness whose prime factors depend on the
curve, and it does not contradict A19's terminating
search for a specified finite support on one specified
curve.  It also does not establish the variable-prime
theta-incidence condition in A18.  An Artin-prime
density theorem continues to supply arithmetic
candidates, while an actual theta incidence at an
appropriate order remains a separate geometric
condition.
